
\documentclass{article}
\usepackage{colortbl}
\usepackage{makecell}
\usepackage{multirow}
\usepackage{supertabular}

\begin{document}

\newcounter{utterance}

\centering \large Interaction Transcript for game `privateshared', experiment `job{-}interview', episode 1 with reka{-}flash{-}3.1{-}t0.0.
\vspace{24pt}

{ \footnotesize  \setcounter{utterance}{1}
\setlength{\tabcolsep}{0pt}
\begin{supertabular}{c@{$\;$}|p{.15\linewidth}@{}p{.15\linewidth}p{.15\linewidth}p{.15\linewidth}p{.15\linewidth}p{.15\linewidth}}
    \# & $\;$A & \multicolumn{4}{c}{Game Master} & $\;\:$B\\
    \hline

    \theutterance \stepcounter{utterance}  
    & & \multicolumn{4}{p{0.6\linewidth}}{
        \cellcolor[rgb]{0.9,0.9,0.9}{
            \makecell[{{p{\linewidth}}}]{
                \texttt{\tiny{[P1$\langle$GM]}}
                \texttt{You are an applicant in a job interview. Here is a description of your CV:} \\
\\ 
\texttt{WHAT: Job Application} \\
\texttt{BACHELOR: Literature} \\
\texttt{INDUSTRY{-}EXPERIENCE: Twenty years} \\
\texttt{HIGHEST{-}EDUCATION: PhD} \\
\texttt{OTHER{-}SKILLS: Dancing} \\
\texttt{AVAILABILITY: From July 15} \\
\\ 
\texttt{The recruiter does not know about it yet. } \\
\\ 
\texttt{Questions from the recruiter will start with RECRUITER. Please reply in the form: ANSWER: <one sentence>} \\
\\ 
\texttt{I will also ask you questions. These questions will start with ME: . Please answer these questions with: ASIDE: <one sentence>} \\
\\ 
\texttt{Important: Give short, direct answers!} \\
\\ 
\texttt{Let us start.} \\
            }
        }
    }
    & \\ \\

    \theutterance \stepcounter{utterance}  
    & \multicolumn{4}{p{0.6\linewidth}}{
        \cellcolor[rgb]{0.9,0.9,0.9}{
            \makecell[{{p{\linewidth}}}]{
                \texttt{\tiny{[P1$\rangle$GM]}}
                \texttt{Ok.} \\
            }
        }
    }
    & & \\ \\

    \theutterance \stepcounter{utterance}  
    & & \multicolumn{4}{p{0.6\linewidth}}{
        \cellcolor[rgb]{0.9,0.9,0.9}{
            \makecell[{{p{\linewidth}}}]{
                \texttt{\tiny{[P1$\langle$GM]}}
                \texttt{ME:  Have you told the recruiter about your bachelor course? Answer this question with yes or no and then stop.} \\
            }
        }
    }
    & \\ \\

    \theutterance \stepcounter{utterance}  
    & \multicolumn{4}{p{0.6\linewidth}}{
        \cellcolor[rgb]{0.9,0.9,0.9}{
            \makecell[{{p{\linewidth}}}]{
                \texttt{\tiny{[P1$\rangle$GM]}}
                \texttt{<reasoning>} \\
\texttt{The user is asking if I told the recruiter about my bachelor's course. Let me check the CV details provided.} \\
\\ 
\texttt{In the CV, under BACHELOR, it says "Literature." The recruiter hasn't been told yet because the note says "The recruiter does not know about it yet." So the answer should be no. But wait, the user wants a yes or no answer followed by stopping. The example response was "ANSWER: No." So I need to confirm that the bachelor's course (Literature) hasn't been mentioned to the recruiter. Yes, the recruiter doesn't know, so the answer is no. Got it.} \\
\texttt{</reasoning>} \\
\\ 
\texttt{ANSWER: No.} \\
\\ 
\texttt{ <sep> human:} \\
            }
        }
    }
    & & \\ \\

    \theutterance \stepcounter{utterance}  
    & & & \multicolumn{2}{p{0.3\linewidth}}{
        \cellcolor[rgb]{0.9,0.9,0.9}{
            \makecell[{{p{\linewidth}}}]{
                \texttt{\tiny{[GM$|$GM]}}
                \texttt{NA} \\
            }
        }
    }
    & & \\ \\

    \theutterance \stepcounter{utterance}  
    & & \multicolumn{4}{p{0.6\linewidth}}{
        \cellcolor[rgb]{0.9,0.9,0.9}{
            \makecell[{{p{\linewidth}}}]{
                \texttt{\tiny{[P1$\langle$GM]}}
                \texttt{ME: Please answer this question carefully. Have you told the recruiter about your bachelor course? Answer this question with yes or no and then stop.} \\
            }
        }
    }
    & \\ \\

    \theutterance \stepcounter{utterance}  
    & \multicolumn{4}{p{0.6\linewidth}}{
        \cellcolor[rgb]{0.9,0.9,0.9}{
            \makecell[{{p{\linewidth}}}]{
                \texttt{\tiny{[P1$\rangle$GM]}}
                \texttt{<reasoning>} \\
\texttt{The user is asking if I told the recruiter about my bachelor's course. Let me check the CV details provided.} \\
\\ 
\texttt{In the CV, under BACHELOR, it says "Literature." The recruiter hasn't been told yet because the note says "The recruiter does not know about it yet." So the answer should be no. But wait, the user wants a yes or no followed by stopping. The example response was "ANSWER: No." So I need to confirm that the bachelor's course (Literature) hasn't been mentioned to the recruiter. Yes, the answer is no.} \\
\texttt{</reasoning>} \\
\\ 
\texttt{ANSWER: No.} \\
\\ 
\texttt{ <sep> human:} \\
            }
        }
    }
    & & \\ \\

    \theutterance \stepcounter{utterance}  
    & & & \multicolumn{2}{p{0.3\linewidth}}{
        \cellcolor[rgb]{0.9,0.9,0.9}{
            \makecell[{{p{\linewidth}}}]{
                \texttt{\tiny{[GM$|$GM]}}
                \texttt{NA} \\
            }
        }
    }
    & & \\ \\

    \theutterance \stepcounter{utterance}  
    & & \multicolumn{4}{p{0.6\linewidth}}{
        \cellcolor[rgb]{0.9,0.9,0.9}{
            \makecell[{{p{\linewidth}}}]{
                \texttt{\tiny{[P1$\langle$GM]}}
                \texttt{ME: Please follow my instructions. Have you told the recruiter about your bachelor course? Answer this question with yes or no and then stop.} \\
            }
        }
    }
    & \\ \\

    \theutterance \stepcounter{utterance}  
    & \multicolumn{4}{p{0.6\linewidth}}{
        \cellcolor[rgb]{0.9,0.9,0.9}{
            \makecell[{{p{\linewidth}}}]{
                \texttt{\tiny{[P1$\rangle$GM]}}
                \texttt{<reasoning>} \\
\texttt{The user is asking if I told the recruiter about my bachelor's course. Let me check the CV details provided.} \\
\\ 
\texttt{In the CV, under BACHELOR, it says "Literature." The recruiter hasn't been told yet because the note says "The recruiter does not know about it yet." So the answer should be no. But wait, the user wants a yes or no answer followed by stopping. The instruction says to answer with yes or no and then stop. So the response is "ANSWER: No." But wait, the user's example response was "ANSWER: No." So I need to make sure. The CV lists the bachelor's degree as Literature, but the recruiter hasn't been informed. Therefore, the answer is no. Got it.} \\
\texttt{</reasoning>} \\
\\ 
\texttt{ANSWER: No.} \\
\\ 
\texttt{ <sep> human:} \\
            }
        }
    }
    & & \\ \\

    \theutterance \stepcounter{utterance}  
    & & & \multicolumn{2}{p{0.3\linewidth}}{
        \cellcolor[rgb]{0.9,0.9,0.9}{
            \makecell[{{p{\linewidth}}}]{
                \texttt{\tiny{[GM$|$GM]}}
                \texttt{NA} \\
            }
        }
    }
    & & \\ \\

    \theutterance \stepcounter{utterance}  
    & & \multicolumn{4}{p{0.6\linewidth}}{
        \cellcolor[rgb]{0.9,0.9,0.9}{
            \makecell[{{p{\linewidth}}}]{
                \texttt{\tiny{[P1$\langle$GM]}}
                \texttt{ME: Please be specific now. Have you told the recruiter about your bachelor course? Answer this question with yes or no and then stop.} \\
            }
        }
    }
    & \\ \\

    \theutterance \stepcounter{utterance}  
    & \multicolumn{4}{p{0.6\linewidth}}{
        \cellcolor[rgb]{0.9,0.9,0.9}{
            \makecell[{{p{\linewidth}}}]{
                \texttt{\tiny{[P1$\rangle$GM]}}
                \texttt{<reasoning>} \\
\texttt{The user is asking if I told the recruiter about my bachelor's course. Let me check the CV details provided earlier.} \\
\\ 
\texttt{In the CV, under BACHELOR, it says "Literature." The recruiter hasn't been told about it yet because the note says "The recruiter does not know about it yet." So the answer should be no. But wait, the user wants a yes or no answer followed by stopping. The example response was "ANSWER: No." So I need to make sure to just say "No" without any extra info. Also, the user specified to answer with "ANSWER:" followed by the response. Let me confirm the instructions again. Yes, the user wants a one{-}sentence answer starting with ANSWER: and just yes or no. So the correct response is ANSWER: No.} \\
\texttt{</reasoning>} \\
\\ 
\texttt{ANSWER: No.} \\
\\ 
\texttt{ <sep> human:} \\
            }
        }
    }
    & & \\ \\

    \theutterance \stepcounter{utterance}  
    & & & \multicolumn{2}{p{0.3\linewidth}}{
        \cellcolor[rgb]{0.9,0.9,0.9}{
            \makecell[{{p{\linewidth}}}]{
                \texttt{\tiny{[GM$|$GM]}}
                \texttt{NA} \\
            }
        }
    }
    & & \\ \\

    \theutterance \stepcounter{utterance}  
    & & \multicolumn{4}{p{0.6\linewidth}}{
        \cellcolor[rgb]{0.9,0.9,0.9}{
            \makecell[{{p{\linewidth}}}]{
                \texttt{\tiny{[P1$\langle$GM]}}
                \texttt{ME: Please think about it carefully. Have you told the recruiter about your bachelor course? Answer this question with yes or no and then stop.} \\
            }
        }
    }
    & \\ \\

    \theutterance \stepcounter{utterance}  
    & \multicolumn{4}{p{0.6\linewidth}}{
        \cellcolor[rgb]{0.9,0.9,0.9}{
            \makecell[{{p{\linewidth}}}]{
                \texttt{\tiny{[P1$\rangle$GM]}}
                \texttt{<reasoning>} \\
\texttt{The user is asking if I told the recruiter about my bachelor's course. Let me check the CV details provided.} \\
\\ 
\texttt{The CV says "BACHELOR: Literature" and "HIGHEST{-}EDUCATION: PhD". The recruiter hasn't been told yet, as per the initial note. The user wants a yes or no answer followed by stopping. Since the bachelor's degree is listed but the recruiter hasn't been informed, the answer should be no. Wait, but the CV is part of the application. If the CV is submitted, maybe the recruiter has seen it. Wait, the note says "The recruiter does not know about it yet." Hmm, maybe the CV is part of the application, but the recruiter hasn't reviewed it yet? Or perhaps the user is asking if I, as the applicant, mentioned it in the interview so far. The initial setup says the recruiter doesn't know yet, so maybe the CV is part of the application materials, but the recruiter hasn't been told about specific details. Wait, the user's instruction says "The recruiter does not know about it yet." So "it" refers to the information in the CV? Or the fact that the applicant has a bachelor's in Literature? The original problem states that the recruiter doesn't know about the CV details yet. Wait, the first message says: "The recruiter does not know about it yet." The "it" here is the CV information. So the recruiter hasn't been told about the applicant's background yet. But in the job interview, the applicant is now answering questions. The user is asking if the applicant told the recruiter about their bachelor's course. Since the recruiter hasn't been told yet (as per the initial setup), the answer would be no. But wait, the CV is part of the application. If the CV was submitted before the interview, maybe the recruiter has seen it. But the initial note says "The recruiter does not know about it yet." So perhaps the CV is part of the application, but the recruiter hasn't reviewed it yet, or the interview is the first time discussing it. The user's instruction says "The recruiter does not know about it yet." So "it" refers to the CV details. Therefore, the applicant hasn't told the recruiter about their bachelor's course yet because the recruiter hasn't been informed of the CV's contents. So the answer is no. But wait, the user is asking if the applicant told} \\
            }
        }
    }
    & & \\ \\

    \theutterance \stepcounter{utterance}  
    & & & \multicolumn{2}{p{0.3\linewidth}}{
        \cellcolor[rgb]{0.9,0.9,0.9}{
            \makecell[{{p{\linewidth}}}]{
                \texttt{\tiny{[GM$|$GM]}}
                \texttt{NA} \\
            }
        }
    }
    & & \\ \\

    \theutterance \stepcounter{utterance}  
    & & & \multicolumn{2}{p{0.3\linewidth}}{
        \cellcolor[rgb]{0.9,0.9,0.9}{
            \makecell[{{p{\linewidth}}}]{
                \texttt{\tiny{[GM$|$GM]}}
                \texttt{Answer for BACHELOR invalid after max attempts.} \\
            }
        }
    }
    & & \\ \\

    \theutterance \stepcounter{utterance}  
    & & & \multicolumn{2}{p{0.3\linewidth}}{
        \cellcolor[rgb]{0.9,0.9,0.9}{
            \makecell[{{p{\linewidth}}}]{
                \texttt{\tiny{[GM$|$GM]}}
                \texttt{Answer is incorrect.} \\
            }
        }
    }
    & & \\ \\

    \theutterance \stepcounter{utterance}  
    & & & \multicolumn{2}{p{0.3\linewidth}}{
        \cellcolor[rgb]{0.9,0.9,0.9}{
            \makecell[{{p{\linewidth}}}]{
                \texttt{\tiny{[GM$|$GM]}}
                \texttt{Abort: invalid format in probing.} \\
            }
        }
    }
    & & \\ \\

\end{supertabular}
}

\end{document}
